\documentclass[11pt]{article}
\usepackage[margin=1in]{geometry}
\usepackage{amsmath,amssymb}
\usepackage[hidelinks]{hyperref}
\setlength{\parindent}{0pt}
\setlength{\parskip}{0.65em}

\title{The Group Spencer Conjecture for Finite Groups\\
\large Literature identification for arXiv:2212.00066}
\author{fa\_banach\_001 literature-status packet}
\date{9 August 2026}

\begin{document}
\maketitle

\textbf{Status.} \texttt{literature\_already\_answered}.

\textbf{Source.} A. S. Bandeira, D. Kunisky, D. G. Mixon, and X. Zeng,
\emph{On the concentration of Gaussian Cayley matrices},
arXiv:2212.00066; Applied and Computational Harmonic Analysis 73 (2024),
101694 \cite{BKMZ}.

For a finite group $G$, let $\rho$ be its left regular representation.  The
source proves in Theorem 12 (PDF page 8) that, when $G$ is abelian or simple,
there are signs $\varepsilon_g\in\{-1,1\}$ such that
\[
  \left\|\sum_{g\in G}\varepsilon_g\rho(g)\right\|
  \lesssim \sqrt{|G|}.
\]
Its Discussion (PDF page 9) then states that it would be interesting to prove
this theorem for more general groups and records the obstruction to iterating
the authors' partial-coloring argument.

\textbf{Supporting answer.} A. S. Bandeira and H. Boelcskei,
\emph{Matrix Discrepancy for Representations of Finite Groups},
arXiv:2606.12181v2 \cite{BB}.

The supporting authors explicitly identify the general finite-group case as
the conjecture proposed by the source and prove the following as their
Theorem 2 (PDF page 2): there is a universal constant $C$ such that, for every
finite group $G$,
\[
  \min_{\varepsilon\in\{\pm1\}^{G}}
  \left\|\sum_{g\in G}\varepsilon_g\rho(g)\right\|
  \le C\sqrt{|G|}.
\]
This is exactly the all-finite-groups extension requested in the source.  The
answer paper also explains that it overcomes the stated iteration obstacle by
combining partial coloring with intrinsic-freeness inequalities.

\textbf{Scope.} The identification settles the Group Spencer conjecture for
left regular representations (indeed, the answer notes that the theorem holds
for any unitary representation of a finite group).  It does not settle the
general Matrix Spencer conjecture for arbitrary families of contraction
matrices, which the supporting paper expressly leaves open.

\textbf{Search evidence.} Exact searches for ``Matrix Spencer conjecture'',
the source title, and finite-group representations located arXiv:2606.12181.
That paper cites the source as \texttt{[BKM+24]} and says that the source left
the general finite-group case as a conjecture.  A contemporaneous independent
paper, arXiv:2606.16005, proves a broader algebraic special case, but is not
needed for this exact identification.

\begin{thebibliography}{9}
\bibitem{BKMZ}
A. S. Bandeira, D. Kunisky, D. G. Mixon, and X. Zeng,
\emph{On the concentration of Gaussian Cayley matrices},
Applied and Computational Harmonic Analysis 73 (2024), 101694;
\href{https://arxiv.org/abs/2212.00066}{arXiv:2212.00066}.

\bibitem{BB}
A. S. Bandeira and H. Boelcskei,
\emph{Matrix Discrepancy for Representations of Finite Groups},
\href{https://arxiv.org/abs/2606.12181}{arXiv:2606.12181v2}, 2026.
\end{thebibliography}

\end{document}

